% Options for packages loaded elsewhere
\PassOptionsToPackage{unicode}{hyperref}
\PassOptionsToPackage{hyphens}{url}
\PassOptionsToPackage{dvipsnames,svgnames,x11names}{xcolor}
\documentclass[
  a4paper,
]{article}
\usepackage{xcolor}
\usepackage[margin=25mm]{geometry}
\usepackage{amsmath,amssymb}
\setcounter{secnumdepth}{-\maxdimen} % remove section numbering
\usepackage{iftex}
\ifPDFTeX
  \usepackage[T1]{fontenc}
  \usepackage[utf8]{inputenc}
  \usepackage{textcomp} % provide euro and other symbols
\else % if luatex or xetex
  \usepackage{unicode-math} % this also loads fontspec
  \defaultfontfeatures{Scale=MatchLowercase}
  \defaultfontfeatures[\rmfamily]{Ligatures=TeX,Scale=1}
\fi
\usepackage{lmodern}
\ifPDFTeX\else
  % xetex/luatex font selection
  \setmainfont[]{DejaVu Serif}
  \setsansfont[]{DejaVu Sans}
  \setmonofont[]{DejaVu Sans Mono}
\fi
% Use upquote if available, for straight quotes in verbatim environments
\IfFileExists{upquote.sty}{\usepackage{upquote}}{}
\IfFileExists{microtype.sty}{% use microtype if available
  \usepackage[]{microtype}
  \UseMicrotypeSet[protrusion]{basicmath} % disable protrusion for tt fonts
}{}
\makeatletter
\@ifundefined{KOMAClassName}{% if non-KOMA class
  \IfFileExists{parskip.sty}{%
    \usepackage{parskip}
  }{% else
    \setlength{\parindent}{0pt}
    \setlength{\parskip}{6pt plus 2pt minus 1pt}}
}{% if KOMA class
  \KOMAoptions{parskip=half}}
\makeatother
\usepackage{longtable,booktabs,array}
\newcounter{none} % for unnumbered tables
\usepackage{calc} % for calculating minipage widths
% Correct order of tables after \paragraph or \subparagraph
\usepackage{etoolbox}
\makeatletter
\patchcmd\longtable{\par}{\if@noskipsec\mbox{}\fi\par}{}{}
\makeatother
% Allow footnotes in longtable head/foot
\IfFileExists{footnotehyper.sty}{\usepackage{footnotehyper}}{\usepackage{footnote}}
\makesavenoteenv{longtable}
\usepackage{graphicx}
\makeatletter
\newsavebox\pandoc@box
\newcommand*\pandocbounded[1]{% scales image to fit in text height/width
  \sbox\pandoc@box{#1}%
  \Gscale@div\@tempa{\textheight}{\dimexpr\ht\pandoc@box+\dp\pandoc@box\relax}%
  \Gscale@div\@tempb{\linewidth}{\wd\pandoc@box}%
  \ifdim\@tempb\p@<\@tempa\p@\let\@tempa\@tempb\fi% select the smaller of both
  \ifdim\@tempa\p@<\p@\scalebox{\@tempa}{\usebox\pandoc@box}%
  \else\usebox{\pandoc@box}%
  \fi%
}
% Set default figure placement to htbp
\def\fps@figure{htbp}
\makeatother
\setlength{\emergencystretch}{3em} % prevent overfull lines
\providecommand{\tightlist}{%
  \setlength{\itemsep}{0pt}\setlength{\parskip}{0pt}}
\usepackage{luatexja}
\usepackage{xurl}
\usepackage{bookmark}
\IfFileExists{xurl.sty}{\usepackage{xurl}}{} % add URL line breaks if available
\urlstyle{same}
\hypersetup{
  pdftitle={The EPR Paradox Is Not a Paradox},
  pdfauthor={Noriaki Kihara (WF System Co., Ltd.) ORCID 0009-0004-6753-4020},
  colorlinks=true,
  linkcolor={Maroon},
  filecolor={Maroon},
  citecolor={Blue},
  urlcolor={Blue},
  pdfcreator={LaTeX via pandoc}}

\title{The EPR Paradox Is Not a Paradox}
\usepackage{etoolbox}
\makeatletter
\providecommand{\subtitle}[1]{% add subtitle to \maketitle
  \apptocmd{\@title}{\par {\large #1 \par}}{}{}
}
\makeatother
\subtitle{― A Supplement to ``Observation as Simultaneous Emission and
Absorption'': Reading the Position--Momentum Pair as a Closure Identity
―}
\author{Noriaki Kihara (WF System Co., Ltd.) ORCID 0009-0004-6753-4020}
\date{November 1, 2026 Version DOI 10.5281/zenodo.23860606}

\ltjsetparameter{jacharrange={-2,-3}}
\begin{document}
\maketitle

\textbf{Author:} Noriaki Kihara (WF System Co., Ltd.)　ORCID
0009-0004-6753-4020\\
\textbf{Date:} 2026-11-01　　\textbf{Version:} v1\\
\textbf{Version DOI:} 10.5281/zenodo.23860606　　\textbf{Concept DOI:}
10.5281/zenodo.23860605\\
\textbf{Note:} The Japanese version
(\texttt{paper\_\allowbreak{}epr\_\allowbreak{}not\_\allowbreak{}a\_\allowbreak{}paradox\_\allowbreak{}supplement\_\allowbreak{}ja\_\allowbreak{}v1.md}) is the
authoritative text; this English version is a translation.

\begin{center}\rule{0.5\linewidth}{0.5pt}\end{center}

\subsection{Abstract}\label{abstract}

The EPR argument {[}X1{]} has three steps. (i) Measuring the position of
one member A of a pair predicts with certainty the position of the other
member B, and measuring the momentum predicts B's momentum (perfect
correlation). (ii) The measurement at A does not disturb the distant B
(locality). (iii) Hence B's position and momentum are simultaneously
elements of reality, and the quantum-mechanical description is
incomplete. This paper takes the readout rule of the preceding paper
``Observation as Simultaneous Emission and Absorption'' {[}S1{]} as its
premise and re-reads the argument. The pair is a closure \(q_A+q_B=0\)
that held once at the source, from which
\(\boldsymbol{x}_B=-\boldsymbol{x}_A\) and
\(\boldsymbol{p}_B=-\boldsymbol{p}_A\) follow simultaneously as
identities. A single end has no value (the one-sided distribution is
uniform), nothing happens at B when A is read, and the correlation
depends neither on the order of the readouts nor on the distance between
the two. The element of reality is the relation called the closure, not
an attribute of one end. Step (ii) is correct, step (i) is a relational
identity, and only the attribution in step (iii) is mistaken. That a
destructive readout exists, drops the chosen variable to one of two
values and severs the relation with the conjugate variable, is placed as
a known fact (resonant absorption, photodetection, homodyne detection
{[}X3{]}); deriving its mechanism is a separate task. ``Instantaneous
selection'' is not written: what can be shown is only that no
finite-speed influence exists {[}X6{]} {[}X7{]} {[}X8{]}. No new number
appears. The claim of this paper is that, given the preceding paper, EPR
is neither a paradox nor a mystery but a misreading of a relational
identity.

\begin{center}\rule{0.5\linewidth}{0.5pt}\end{center}

\subsection{1. The EPR argument}\label{the-epr-argument}

In 1935 Einstein, Podolsky and Rosen set up the following argument
{[}X1{]}. Two particles are in a state satisfying \(x_1-x_2=x_0\) and
\(p_1+p_2=0\). Measuring the position of particle 1 predicts the
position of particle 2 with certainty, and measuring its momentum
predicts the momentum of particle 2, without touching particle 2. If the
criterion of reality is ``a quantity that can be predicted with
certainty without disturbing the system has a corresponding element of
reality,'' then both the position and the momentum of particle 2 are
elements of reality. Since quantum mechanics cannot describe both at
once, it is incomplete.

Bohr's reply {[}X2{]} was that the measurement on particle 1 does not
mechanically disturb particle 2, but defines the conditions under which
something about particle 2 can be predicted. The question later moved,
through Bell's theorem {[}X18{]}, to the exclusion of local hidden
variables, and in continuous variables the original EPR form was
realized through Reid's criterion {[}X4{]} and the experiment of Ou et
al.~{[}X3{]}. A review is given by Reid et al.~{[}X5{]}.

This paper decides which part of the argument is right and which is
wrong using only the three premises of the preceding paper {[}S1{]}. No
new premise is added. Only the existence of a destructive readout is
brought in from outside, as a known fact.

\begin{center}\rule{0.5\linewidth}{0.5pt}\end{center}

\subsection{2. Premises}\label{premises}

\subsubsection{2.1 The three premises of the preceding
paper}\label{the-three-premises-of-the-preceding-paper}

The preceding paper {[}S1{]} placed the following three premises and
derived the form of the Born rule, the two-party correlation
\(E=\mp\cos n\theta\), the CHSH bound \(2\sqrt2\), and no-signaling.

\begin{itemize}
\tightlist
\item
  (P1) The distance between two parties is set by the ratio of a shared
  viewing angle to the partner's size; only ratios are readable. The
  overall orientation and size (the gauge) are not readable.
\item
  (P2) An interaction is a single relation in which emission and
  absorption hold simultaneously; there is no before and after.
\item
  (P3) The phase of a relation carries a harmonic order \(n\) fixed by
  the kind of wave.
\end{itemize}

What this paper uses is the gauge of P1 (the absolute direction is
unreadable from one side alone) and P2 (a single relation has no before
and after).

\subsubsection{2.2 The pair is a closure}\label{the-pair-is-a-closure}

Two states \(q_A, q_B\) form a pair from one source when, at the source,

\[
q_A+q_B=0
\]

holds {[}S3{]}. This is not a condition; it is the identity expressing
that the source tied the two ends as one relation. From then on
\(q_B=-q_A\) holds at all times.

\subsubsection{2.3 Reading position and
momentum}\label{reading-position-and-momentum}

In the readout plane of the preceding paper, the position vector
\(\boldsymbol{x}\) of a state \(q\) is read as the image of \(q\) in the
readout plane, and the momentum vector \(\boldsymbol{p}\) as a quantity
proportional to its rate of change in time. In this series the position
is a relational quantity read as the quotient angle of two states
{[}S1{]} {[}S2{]}; no coordinates are brought in from outside. The
source is taken as the origin.

\subsubsection{2.4 Destructive readouts
exist}\label{destructive-readouts-exist}

The ``first interaction'' of §1 of the preceding paper, the one that
switches the partner's discrete state, is known. Resonant absorption,
the absorption of a photon by a photodetector, and, for continuous
variables, homodyne detection are of this kind. Homodyne detection
interferes the field with a local oscillator and reads either the \(x\)
quadrature or the \(p\) quadrature, absorbing the light {[}X3{]}
{[}X5{]}. This paper places it as a fact.

\begin{itemize}
\tightlist
\item
  (P4, a known fact) A destructive readout exists. It drops the chosen
  variable to one of two values and severs, for the same pair, the
  relation with the conjugate variable.
\end{itemize}

The mechanism of P4, i.e.~deriving this interaction from the dynamics of
the closure, is not the subject of this paper (§6).

\begin{center}\rule{0.5\linewidth}{0.5pt}\end{center}

\subsection{3. Identities}\label{identities}

\subsubsection{3.1 One closure fixes position and momentum at
once}\label{one-closure-fixes-position-and-momentum-at-once}

Since \(q_B=-q_A\) holds at all times, for the image in the readout
plane and its rate of change,

\[
\boldsymbol{x}_B=-\boldsymbol{x}_A,\qquad \boldsymbol{p}_B=-\boldsymbol{p}_A
\]

hold simultaneously, by linearity. EPR's \(x_1-x_2=x_0\) and
\(p_1+p_2=0\) are the same equations with the sign convention of one
coordinate changed. That the sum of positions and the sum of momenta are
zero at the same time is a restatement of the closure having held once
at the source, not something that comes to hold after a measurement.

\begin{figure}
\centering
\pandocbounded{\includegraphics[keepaspectratio,alt={The closed pair. (a) In the readout plane with the source S at the origin, q\_B=-q\_A lies at the antipode, and the position vectors \textbackslash boldsymbol\{x\} and the velocity vectors \textbackslash boldsymbol\{v\} are opposite at the same time. (b) As time runs the two stay antipodal. (c) When A is read destructively against the axis \textbackslash boldsymbol\{n\}\_A, A's record falls to one of ±, while the state of B does not change at all. All points are computed from the formulas (図表/make\_figures.py).}]{figs/fig1_closure_pair_en.png}}
\caption{The closed pair. (a) In the readout plane with the source S at
the origin, \(q_B=-q_A\) lies at the antipode, and the position vectors
\(\boldsymbol{x}\) and the velocity vectors \(\boldsymbol{v}\) are
opposite at the same time. (b) As time runs the two stay antipodal. (c)
When A is read destructively against the axis \(\boldsymbol{n}_A\), A's
record falls to one of ±, while the state of B does not change at all.
All points are computed from the formulas
(\protect\texttt{図表/\allowbreak{}make\_\allowbreak{}figures.py}).}
\end{figure}

\subsubsection{3.2 A single end has no
value}\label{a-single-end-has-no-value}

The overall phase (the orientation of the pair relative to the source)
is the gauge of P1 and is unreadable. Reading the sign against an axis
\(\boldsymbol{n}_A\) on one side alone, the distribution averaged over
the gauge is uniform, \(\tfrac12\) for two values (3.5 of the preceding
paper). Neither ``the position of B'' nor ``the momentum of B'' exists
as a value belonging to B alone. What exists is the closure, i.e.~the
relations \(\boldsymbol{x}_A+\boldsymbol{x}_B=0\) and
\(\boldsymbol{p}_A+\boldsymbol{p}_B=0\).

\subsubsection{3.3 The correlation depends neither on order nor on
distance}\label{the-correlation-depends-neither-on-order-nor-on-distance}

The cross term of the pair is

\[
\mathrm{Sc}\,[q_A\bar q_B]=-|q_A|^{2},
\]

independent of the overall phase and of time. Whether A is read first, B
is read first, or both are read at once, the result of comparing the
records is the same. It does not depend on the viewing angle (distance)
between the two either.

\begin{figure}
\centering
\pandocbounded{\includegraphics[keepaspectratio,alt={Order independence. The relations from the source S to A and to B are one relation of zero spacetime interval (P2), and the readouts at A and at B are two spacelike-separated events. In one frame A comes first, in another B comes first, but the comparison of the records is the same. The points are computed by Lorentz transformation.}]{figs/fig2_order_independence_en.png}}
\caption{Order independence. The relations from the source S to A and to
B are one relation of zero spacetime interval (P2), and the readouts at
A and at B are two spacelike-separated events. In one frame A comes
first, in another B comes first, but the comparison of the records is
the same. The points are computed by Lorentz transformation.}
\end{figure}

\subsubsection{3.4 B's marginal distribution does not
change}\label{bs-marginal-distribution-does-not-change}

An operation on A acts on \(q_A\) only, not on \(q_B\). The one-sided
distribution of B remains uniform as in 3.2. There is no way to tell on
B's side alone whether A has been read or not (no-signaling, 3.5 of the
preceding paper).

\subsubsection{3.5 The relation with the conjugate variable is
severed}\label{the-relation-with-the-conjugate-variable-is-severed}

When \(x\) is read destructively at A, by P4 the relation of \(p\) for
the same pair is severed. This is a known property of P4, not a
derivation of this paper. EPR's point that one may freely choose to
measure position or momentum stands as it is; only the correlation of
the chosen variable remains.

\begin{center}\rule{0.5\linewidth}{0.5pt}\end{center}

\subsection{4. Re-reading}\label{re-reading}

The three steps of EPR are re-read with §3.

\begin{itemize}
\tightlist
\item
  \begin{enumerate}
  \def\labelenumi{(\roman{enumi})}
  \tightlist
  \item
    The perfect correlation is correct. But it is not a consequence of
    measurement; it is the closure identity 3.1. It holds before
    anything is measured.
  \end{enumerate}
\item
  \begin{enumerate}
  \def\labelenumi{(\roman{enumi})}
  \setcounter{enumi}{1}
  \tightlist
  \item
    Locality is correct. Nothing happens at B when A is read (3.4).
  \end{enumerate}
\item
  \begin{enumerate}
  \def\labelenumi{(\roman{enumi})}
  \setcounter{enumi}{2}
  \tightlist
  \item
    Elements of reality exist. But they are not one-sided attributes
    such as ``the position of B'' or ``the momentum of B''; they are the
    relations \(\boldsymbol{x}_A+\boldsymbol{x}_B=0\) and
    \(\boldsymbol{p}_A+\boldsymbol{p}_B=0\). A single end has no value
    (3.2). What quantum mechanics describes is this relation, and as a
    description of the relation it is complete. It looked incomplete
    because a relational quantity was attributed to one end.
  \end{enumerate}
\end{itemize}

Hence there is no paradox. Nor is the picture needed that, when the
position is read at A, ``the opposite selection occurs instantaneously
at B.'' There is no event at B; one end of what the closure had already
fixed was read at A.

A word on ``instantaneous.'' The time order of two spacelike-separated
events depends on the frame, so the sentence ``the readout at A
instantaneously fixed the value at B'' has no frame-independent meaning.
What experiments can show is a lower bound: if there were a finite-speed
influence, its speed would exceed \(10^4\) times the speed of light
{[}X6{]} {[}X7{]}; moreover, assuming a finite-speed influence allows
signaling {[}X8{]}. ``No finite-speed selection exists'' is the most
that can be shown, and the reading of this paper is consistent with it.

A remark on comparing the axes of the two detectors. Since \(E\) is the
inner product \(\boldsymbol{n}_A\cdot\boldsymbol{n}_B\), no common
coordinates at the two ends are required; but in curved spacetime or
over long distances the comparison depends on the path along which an
axis is transported, and the correlation is reduced by precession
{[}X9{]}. This is a manifestation of P1's ``ratios close among three
bodies'' depending on which third party mediates.

\begin{center}\rule{0.5\linewidth}{0.5pt}\end{center}

\subsection{5. Relation to earlier work}\label{relation-to-earlier-work}

The conclusion of this paper is not new. The claim that a single end has
no value and only correlations are real is in Peres {[}X10{]} and Mermin
{[}X11{]}, and the treatment of EPR in relational quantum mechanics
{[}X12{]} argues that there is no non-local event on B's side. The
locality claim from consistent histories {[}X13{]}, the distinction
between no-signaling and outcome correlation {[}X14{]}, the toy model in
which EPR-type perfect correlations arise locally from an epistemic
restriction alone {[}X15{]}, the review of formulations that fix both
ends at once by an action principle {[}X16{]}, and the transactional
interpretation that depicts quantum transitions as handshakes {[}X17{]}
all point in the same direction.

What belongs to this paper alone is limited to four points. First, the
position is placed as a relational quantity (the quotient angle) without
a background spacetime, and the EPR pair is written as the algebraic
identity \(q_A+q_B=0\) of a closure: a formula, not a stance. Second,
that a single end has no value is derived from the gauge of P1, not
posited as a stance. Third, that the same closure fixes the two
relations of position and momentum at once (3.1) is checked within the
readout of the series. Fourth, it is continuous with the preceding paper
{[}S1{]}, where the Bell side (\(2\sqrt2\)) follows from the bilinearity
of the same readout rule. Since the original EPR argument is built from
perfect correlations alone, this paper treats only the
perfect-correlation side and leaves the non-classicality to the
preceding paper.

\begin{center}\rule{0.5\linewidth}{0.5pt}\end{center}

\subsection{6. How it would fail, and the separate
task}\label{how-it-would-fail-and-the-separate-task}

This paper would fail in three ways.

\begin{itemize}
\tightlist
\item
  If, within the readout of the series, what fixes the two relations of
  position and momentum at once were shown not to be a single closure
  (the negation of 3.1).
\item
  If an experimental fact appeared in which the one-sided distribution
  is not uniform. That would mean signaling is possible, and not only
  this paper but the preceding one would fall.
\item
  If, when the P4 interaction is made concrete from the dynamics of the
  closure, it did not drop to one of two values, or the share were not
  \(\cos^2\). P4 would be false, and this paper and 3.3 of the preceding
  paper would fall.
\end{itemize}

The third is the separate task. The mechanism of the destructive
readout, i.e.~deriving the pull-in of the axis and the switching of the
state from an action, is contained in neither this paper nor the
preceding one. This paper does not depend on that mechanism; it uses
only its existence.

\begin{center}\rule{0.5\linewidth}{0.5pt}\end{center}

\subsection{7. Classification}\label{classification}

{\def\LTcaptype{none} % do not increment counter
\begin{longtable}[]{@{}
  >{\raggedright\arraybackslash}p{(\linewidth - 2\tabcolsep) * \real{0.5000}}
  >{\raggedright\arraybackslash}p{(\linewidth - 2\tabcolsep) * \real{0.5000}}@{}}
\toprule\noalign{}
\begin{minipage}[b]{\linewidth}\raggedright
Class
\end{minipage} & \begin{minipage}[b]{\linewidth}\raggedright
Content
\end{minipage} \\
\midrule\noalign{}
\endhead
\bottomrule\noalign{}
\endlastfoot
Identities & 3.1 from the closure, the sums of positions and of momenta
vanish at once; 3.3 independence of the cross term from order and
distance; 3.4 invariance of B's marginal \\
Readings & 2.3 position = image in the readout plane, momentum = its
rate of change; 3.2 a single end has no value; §4 element of reality =
relation \\
Known observations & P4 the existence of destructive readouts (resonant
absorption, photodetection, homodyne detection {[}X3{]}); the lower
bounds on finite-speed influences {[}X6{]} {[}X7{]} {[}X8{]} \\
How it would fail & the three of §6 \\
\end{longtable}
}

\begin{center}\rule{0.5\linewidth}{0.5pt}\end{center}

\subsection{References}\label{references}

\subsubsection{Self-citations (Concept DOI,
10.5281/zenodo.〔number〕)}\label{self-citations-concept-doi-10.5281zenodo.number}

{[}S1{]} N. Kihara, ``観測は放射と吸収の同時成立である ―
視野角と大きさの読出し規則からボルン則・ベル相関・ツィレルソン限界を導く
―'' (Observation as Simultaneous Emission and Absorption: Deriving the
Born Rule, Bell Correlations and the Tsirelson Bound from a
Viewing-Angle Readout Rule), 2026-11-01. Concept DOI
10.5281/zenodo.23857884.

{[}S2{]} N. Kihara, ``複素数を正しく理解していたか ―
振幅はどこに住むか、虚数読み、\(\mathbb{C}\) → \(\mathbb{H}\) →
\(\mathbb{O}\) の帰属 ―'' (Did We Understand Complex Numbers Correctly?
Where the Amplitude Lives, the Imaginary Reading, and Attribution across
\(\mathbb{C}\) → \(\mathbb{H}\) → \(\mathbb{O}\)), 2026-10-23. Concept
DOI 10.5281/zenodo.23585436.

{[}S3{]} N. Kihara, ``無名等振幅複合波モデル基本公理系 v9'' (Anonymous
Equal-Amplitude Composite-Wave Model: Basic Axiom System v9),
2026-07-24. Concept DOI 10.5281/zenodo.21315735.

\subsubsection{External references}\label{external-references}

{[}X1{]} A. Einstein, B. Podolsky, N. Rosen, ``Can Quantum-Mechanical
Description of Physical Reality Be Considered Complete?,'' Phys.
Rev.~47, 777--780 (1935). DOI 10.1103/PhysRev.47.777

{[}X2{]} N. Bohr, ``Can Quantum-Mechanical Description of Physical
Reality be Considered Complete?,'' Phys. Rev.~48, 696--702 (1935). DOI
10.1103/PhysRev.48.696

{[}X3{]} Z. Y. Ou, S. F. Pereira, H. J. Kimble, K. C. Peng,
``Realization of the Einstein-Podolsky-Rosen paradox for continuous
variables,'' Phys. Rev.~Lett. 68, 3663--3666 (1992). DOI
10.1103/PhysRevLett.68.3663

{[}X4{]} M. D. Reid, ``Demonstration of the Einstein-Podolsky-Rosen
paradox using nondegenerate parametric amplification,'' Phys. Rev.~A 40,
913--923 (1989). DOI 10.1103/PhysRevA.40.913

{[}X5{]} M. D. Reid, P. D. Drummond, W. P. Bowen, E. G. Cavalcanti, P.
K. Lam, H. A. Bachor, U. L. Andersen, G. Leuchs, ``Colloquium: The
Einstein-Podolsky-Rosen paradox: From concepts to applications,''
Rev.~Mod. Phys. 81, 1727--1751 (2009). DOI 10.1103/RevModPhys.81.1727

{[}X6{]} D. Salart, A. Baas, C. Branciard, N. Gisin, H. Zbinden,
``Testing the speed of `spooky action at a distance','' Nature 454,
861--864 (2008). DOI 10.1038/nature07121

{[}X7{]} J. Yin, Y. Cao, H.-L. Yong, J.-G. Ren, H. Liang, S.-K. Liao, F.
Zhou, C. Liu, Y.-P. Wu, G.-S. Pan, L. Li, N.-L. Liu, Q. Zhang, C.-Z.
Peng, J.-W. Pan, ``Lower Bound on the Speed of Nonlocal Correlations
without Locality and Measurement Choice Loopholes,'' Phys. Rev.~Lett.
110, 260407 (2013). DOI 10.1103/PhysRevLett.110.260407

{[}X8{]} J.-D. Bancal, S. Pironio, A. Acín, Y.-C. Liang, V. Scarani, N.
Gisin, ``Quantum non-locality based on finite-speed causal influences
leads to superluminal signalling,'' Nat. Phys. 8, 867--870 (2012). DOI
10.1038/nphys2460

{[}X9{]} H. Terashima, M. Ueda, ``Einstein-Podolsky-Rosen correlation in
a gravitational field,'' Phys. Rev.~A 69, 032113 (2004). DOI
10.1103/PhysRevA.69.032113

{[}X10{]} A. Peres, ``Unperformed experiments have no results,'' Am. J.
Phys. 46, 745--747 (1978). DOI 10.1119/1.11393

{[}X11{]} N. D. Mermin, ``What is quantum mechanics trying to tell
us?,'' Am. J. Phys. 66, 753--767 (1998). DOI 10.1119/1.18955

{[}X12{]} M. Smerlak, C. Rovelli, ``Relational EPR,'' Found. Phys. 37,
427--445 (2007). DOI 10.1007/s10701-007-9105-0

{[}X13{]} R. B. Griffiths, ``EPR, Bell, and quantum locality,'' Am. J.
Phys. 79, 954--965 (2011). DOI 10.1119/1.3606371

{[}X14{]} J. P. Jarrett, ``On the Physical Significance of the Locality
Conditions in the Bell Arguments,'' Noûs 18, 569--589 (1984). DOI
10.2307/2214878

{[}X15{]} R. W. Spekkens, ``Evidence for the epistemic view of quantum
states: A toy theory,'' Phys. Rev.~A 75, 032110 (2007). DOI
10.1103/PhysRevA.75.032110

{[}X16{]} K. B. Wharton, N. Argaman, ``Colloquium: Bell's theorem and
locally mediated reformulations of quantum mechanics,'' Rev.~Mod. Phys.
92, 021002 (2020). DOI 10.1103/RevModPhys.92.021002

{[}X17{]} J. G. Cramer, ``The transactional interpretation of quantum
mechanics,'' Rev.~Mod. Phys. 58, 647--687 (1986). DOI
10.1103/RevModPhys.58.647

{[}X18{]} J. S. Bell, ``On the Einstein Podolsky Rosen paradox,''
Physics Physique Fizika 1, 195--200 (1964). DOI
10.1103/PhysicsPhysiqueFizika.1.195

\begin{center}\rule{0.5\linewidth}{0.5pt}\end{center}

\subsection{Appendix A　Reproducing the figures and the
identities}\label{appendix-a-reproducing-the-figures-and-the-identities}

Figures 1 and 2 are generated by \texttt{図表/\allowbreak{}make\_\allowbreak{}figures.py} in the
same folder. \texttt{python3\ make\_\allowbreak{}figures.py} produces the Japanese
labels and \texttt{python3\ make\_\allowbreak{}figures.py\ en} the English-label
figures (\texttt{fig*\_\allowbreak{}en}), written to \texttt{図表/\allowbreak{}figures/\allowbreak{}} as PNG,
SVG and PDF. All points are computed directly from the formulas; no
random numbers are used. The identities are checked on a fixed grid by
\texttt{図表/\allowbreak{}verify\_\allowbreak{}identities.py} (C1: from \(q_B=-q_A\) the sums of
positions and of velocities vanish at once; C2: the gauge average of the
one-sided sign is \(\tfrac12\); C3: the cross term is independent of the
overall phase and of time; C4: an operation on A does not change B's
distribution). The results are in \texttt{図表/\allowbreak{}results/\allowbreak{}checks.txt}.
Dependencies: numpy and matplotlib.

\subsection{Version history}\label{version-history}

\begin{itemize}
\tightlist
\item
  v1 (2026-11-01): first draft. Written as a supplement right after the
  publication of the preceding paper {[}S1{]} (Version DOI
  10.5281/zenodo.23857885), from the dialogue between the author and
  Claude Code (EPR in position and momentum rather than spin; the
  destructive readout placed as a known fact; ``instantaneous
  selection'' not written; the earlier work placed up front). Published
  on Zenodo (Version DOI 10.5281/zenodo.23860606, Concept DOI
  10.5281/zenodo.23860605; this paper is registered as isSupplementTo
  {[}S1{]}).
\end{itemize}

\end{document}
